% GENERATED FILE. Edit canonical lessons, then run npm run generate:books.
% canonical-source-hash: fnv1a64:19c0aeb1a858d3c2
% canonical-lessons: LA-C369-hortor, LA-C369-dissuadeo, LA-C369-persuasio, LA-C369-excuso, LA-C369-excusatio

\chapter{Urging and Dissuading}
\label{ch:la-urging-and-dissuading}
\hlchaptermodality{369}

\begin{chapteropening}
\textbf{By the end of this chapter:} \emph{I can urge someone on, talk someone out of a plan, speak of persuasion, and offer or turn down an excuse.}

Complete the last lesson of chapter 369: I can urge someone on, talk someone out of a plan, speak of persuasion, and offer or turn down an excuse.
\end{chapteropening}

\section[hortor, hortārī]{hortor, hortārī — to urge, to encourage}

\label{lesson:LA-C369-hortor}



\begin{quote}
Before the new one: say the Latin for to discuss, then the Latin for to speak out.
\end{quote}

\subsection*{You'll want to know: hortor, hortārī}
\textbf{hortor, hortārī} — \textquotedblleft{}to urge, to encourage\textquotedblright{}.

\begin{quote}
Mīlitēs hortātus est ut fortiter pugnārent.
\end{quote}

\begin{quote}
\emph{He urged the soldiers to fight bravely.}
\end{quote}

Deponent: \emph{hortor, hortārī, hortātus sum} + accusative and \emph{ut} (or \emph{nē}) with the subjunctive. English \emph{exhort}.

\begin{cousinweb}
\emph{hortor} comes from an old verb \emph{horior} \textquotedblleft{}to urge on\textquotedblright{} (EtymDB; Lewis \& Short), traced to a root meaning \textquotedblleft{}to desire\textquotedblright{} that Greek \emph{khairō} \textquotedblleft{}to rejoice\textquotedblright{} shares. English \emph{exhort} comes from the compound \emph{exhortor}.
\end{cousinweb}

\begin{grammarlens}[title={What you've built}]
One more word for telling, explaining and describing at B1.
\end{grammarlens}

\subsection*{Your turn}
\begin{itemize}
\raggedright
  \item \emph{Say it:} \emph{hortor, hortārī}
  \item \emph{Say it:} \emph{hortor, hortārī}, once more
  \item \emph{Say it:} a sentence of your own with \emph{hortor, hortārī}
  \item \emph{Recall:} say the Latin for to discuss, then the Latin for to speak out, then say \emph{hortor, hortārī} again
\end{itemize}

\begin{tcolorbox}[breakable,colback=teal!4,colframe=teal!35!black,title={Before you move on}]
What does hortor, hortārī mean? (\textquotedblleft{}To urge, to encourage\textquotedblright{}.) Say it once more.
\end{tcolorbox}
\section[dissuādeō, dissuādēre]{dissuādeō, dissuādēre — to advise against, to dissuade}

\label{lesson:LA-C369-dissuadeo}



\begin{quote}
Before the new one: say the Latin for to speak out, then the Latin for to urge.
\end{quote}

\subsection*{You'll want to know: dissuādeō, dissuādēre}
\textbf{dissuādeō, dissuādēre} — \textquotedblleft{}to advise against, to dissuade\textquotedblright{}.

\begin{quote}
Amīcī mihi hoc iter dissuāsērunt.
\end{quote}

\begin{quote}
\emph{My friends advised me against this journey.}
\end{quote}

\emph{dissuādeō, dissuādēre, dissuāsī, dissuāsum}, with the dative of the person: the opposite of \emph{suādeō}. English \emph{dissuade}.

\begin{cousinweb}
\emph{dissuādeō} is built from \emph{dis-} \textquotedblleft{}apart, away\textquotedblright{} and \emph{suādeō} \textquotedblleft{}to advise\textquotedblright{} (EtymDB; Lewis \& Short). English \emph{dissuade} comes from it.
\end{cousinweb}

\begin{grammarlens}[title={What you've built}]
One more word for telling, explaining and describing at B1.
\end{grammarlens}

\subsection*{Your turn}
\begin{itemize}
\raggedright
  \item \emph{Say it:} \emph{dissuādeō, dissuādēre}
  \item \emph{Say it:} \emph{dissuādeō, dissuādēre}, once more
  \item \emph{Say it:} a sentence of your own with \emph{dissuādeō, dissuādēre}
  \item \emph{Recall:} say the Latin for to speak out, then the Latin for to urge, then say \emph{dissuādeō, dissuādēre} again
\end{itemize}

\begin{tcolorbox}[breakable,colback=teal!4,colframe=teal!35!black,title={Before you move on}]
What does dissuādeō, dissuādēre mean? (\textquotedblleft{}To advise against, to dissuade\textquotedblright{}.) Say it once more.
\end{tcolorbox}
\section[persuāsiō, persuāsiōnis]{persuāsiō, persuāsiōnis — persuasion; a conviction}

\label{lesson:LA-C369-persuasio}



\begin{quote}
Before the new one: say the Latin for to urge, then the Latin for to advise against.
\end{quote}

\subsection*{You'll want to know: persuāsiō, persuāsiōnis}
\textbf{persuāsiō, persuāsiōnis} — \textquotedblleft{}persuasion; a conviction\textquotedblright{}.

\begin{quote}
Persuāsiōne, nōn vī, vincere mālō.
\end{quote}

\begin{quote}
\emph{I prefer to win by persuasion, not by force.}
\end{quote}

\emph{persuāsiō, persuāsiōnis} f., the noun of \emph{persuādeō}. English \emph{persuasion}.

\begin{cousinweb}
\emph{persuāsiō} is the noun built on \emph{persuādeō} \textquotedblleft{}to persuade\textquotedblright{} (EtymDB; Lewis \& Short), itself \emph{per-} \textquotedblleft{}thoroughly\textquotedblright{} with \emph{suādeō} \textquotedblleft{}to advise\textquotedblright{}. English \emph{persuasion} comes from it.
\end{cousinweb}

\begin{grammarlens}[title={What you've built}]
One more word for telling, explaining and describing at B1.
\end{grammarlens}

\subsection*{Your turn}
\begin{itemize}
\raggedright
  \item \emph{Say it:} \emph{persuāsiō, persuāsiōnis}
  \item \emph{Say it:} \emph{persuāsiō, persuāsiōnis}, once more
  \item \emph{Say it:} a sentence of your own with \emph{persuāsiō, persuāsiōnis}
  \item \emph{Recall:} say the Latin for to urge, then the Latin for to advise against, then say \emph{persuāsiō, persuāsiōnis} again
\end{itemize}

\begin{tcolorbox}[breakable,colback=teal!4,colframe=teal!35!black,title={Before you move on}]
What does persuāsiō, persuāsiōnis mean? (\textquotedblleft{}Persuasion; a conviction\textquotedblright{}.) Say it once more.
\end{tcolorbox}
\section[excūsō, excūsāre]{excūsō, excūsāre — to excuse; to plead as an excuse}

\label{lesson:LA-C369-excuso}



\begin{quote}
Before the new one: say the Latin for to advise against, then the Latin for persuasion.
\end{quote}

\subsection*{You'll want to know: excūsō, excūsāre}
\textbf{excūsō, excūsāre} — \textquotedblleft{}to excuse; to plead as an excuse\textquotedblright{}.

\begin{quote}
Morbum excūsāvit.
\end{quote}

\begin{quote}
\emph{He pleaded illness as his excuse.}
\end{quote}

\emph{excūsō, excūsāre, excūsāvī, excūsātum}. \emph{Sē excūsāre}, to excuse oneself. Built on \emph{causa}: out of the case. English \emph{excuse}.

\begin{cousinweb}
\emph{excūsō} is built from \emph{ex} \textquotedblleft{}out of\textquotedblright{} and \emph{causa} \textquotedblleft{}a case, a charge\textquotedblright{} (EtymDB; Lewis \& Short): to take someone out of the charge. English \emph{excuse} comes from it.
\end{cousinweb}

\begin{grammarlens}[title={What you've built}]
One more word for telling, explaining and describing at B1.
\end{grammarlens}

\subsection*{Your turn}
\begin{itemize}
\raggedright
  \item \emph{Say it:} \emph{excūsō, excūsāre}
  \item \emph{Say it:} \emph{excūsō, excūsāre}, once more
  \item \emph{Say it:} a sentence of your own with \emph{excūsō, excūsāre}
  \item \emph{Recall:} say the Latin for to advise against, then the Latin for persuasion, then say \emph{excūsō, excūsāre} again
\end{itemize}

\begin{tcolorbox}[breakable,colback=teal!4,colframe=teal!35!black,title={Before you move on}]
What does excūsō, excūsāre mean? (\textquotedblleft{}To excuse; to plead as an excuse\textquotedblright{}.) Say it once more.
\end{tcolorbox}
\section[excūsātiō, excūsātiōnis]{excūsātiō, excūsātiōnis — an excuse, a plea}

\label{lesson:LA-C369-excusatio}



\begin{quote}
Before the new one: say the Latin for persuasion, then the Latin for to excuse.
\end{quote}

\subsection*{You'll want to know: excūsātiō, excūsātiōnis}
\textbf{excūsātiō, excūsātiōnis} — \textquotedblleft{}an excuse, a plea\textquotedblright{}.

\begin{quote}
Haec excūsātiō nōn accipitur.
\end{quote}

\begin{quote}
\emph{This excuse is not accepted.}
\end{quote}

\emph{excūsātiō, excūsātiōnis} f., the noun of \emph{excūsō}.

\begin{cousinweb}
\emph{excūsātiō} is the noun built on \emph{excūsō} \textquotedblleft{}to excuse\textquotedblright{} (EtymDB; Lewis \& Short), from \emph{ex} and \emph{causa}.
\end{cousinweb}

\begin{grammarlens}[title={What you've built}]
One more word for telling, explaining and describing at B1.
\end{grammarlens}

\subsection*{Your turn}
\begin{itemize}
\raggedright
  \item \emph{Say it:} \emph{excūsātiō, excūsātiōnis}
  \item \emph{Say it:} \emph{excūsātiō, excūsātiōnis}, once more
  \item \emph{Say it:} a sentence of your own with \emph{excūsātiō, excūsātiōnis}
  \item \emph{Recall:} say the Latin for persuasion, then the Latin for to excuse, then say \emph{excūsātiō, excūsātiōnis} again
\end{itemize}

\begin{tcolorbox}[breakable,colback=teal!4,colframe=teal!35!black,title={Before you move on}]
What does excūsātiō, excūsātiōnis mean? (\textquotedblleft{}An excuse, a plea\textquotedblright{}.) Say it once more.
\end{tcolorbox}
